% Vietnamese translation for OI Public Library
% Source-File: IOI中国国家候选队论文/2008/Day1/3.刘弈《浅谈信息学中状态的合理设计与应用》/浅谈信息学中状态的合理设计与应用.ppt
\documentclass[11pt]{article}
\usepackage{oipl-vn}

\newcommand{\Oh}{\mathcal{O}}

\title{Bản trình chiếu: Thiết kế trạng thái hợp lý}
\author{Liu Yi}
\date{2008}

\begin{document}
\maketitle

\origtitle{Slide companion for reasonable state design and applications}
\sourcepath{IOI China candidate team papers / 2008 / Day 1 / Liu Yi.ppt}

\section{Phạm vi bản dịch}

Tệp PowerPoint đi cùng bài viết đã được dịch từ tài liệu Word. Bản ghi chú này
giữ lại trọng tâm: số lượng trạng thái quan trọng không kém chi phí xử lý từng
trạng thái.

\section{Công thức tư duy}

Trong nhiều thuật toán, thời gian chạy có thể nhìn như
\[
  \text{số trạng thái}\times\text{chi phí xử lý mỗi trạng thái}.
\]
Tối ưu chuyển trạng thái chỉ giải quyết vế thứ hai. Nếu cách biểu diễn trạng
thái chứa quá nhiều thông tin dư, thuật toán vẫn có thể quá chậm dù mỗi bước
xử lý đã được viết tốt.

\section{Ba hướng cải tiến}

Các slide lần lượt nhắc ba hướng: đếm đúng số trạng thái thật sự cần xét,
thay đổi cách biểu diễn trạng thái để gộp nhiều trạng thái tương đương, và
thiết kế lại trạng thái khi mô hình quen thuộc dẫn tới bế tắc. Ví dụ
\textit{Square Root} cho thấy nhóm truy vấn theo modulo nguyên tố có thể biến
một cách liệt kê tưởng như thô thành lời giải đủ nhanh.

\section{Kết luận}

Một trạng thái tốt cần giữ đúng phần thông tin ảnh hưởng đến tương lai và bỏ
phần chỉ là lịch sử hoặc cách biểu diễn. Khi gặp DP, tìm kiếm hoặc liệt kê
quá lớn, câu hỏi đầu tiên nên là: hai trạng thái nào thật sự khác nhau với
phần còn lại của bài toán?

\end{document}
